%% kapa: exam=1 round=34 year=2023 period=2 form=A subject=accounting subject_ko=회계학 seq=37 number=77
%% source: Q-Net 감정평가사 기출문제 (cst00309) artlSeq=5211789 file=raw/1008/5211789/2023년도 제34회 감정평가사 1차시험 2교시.pdf
%% conversion: kordoc 4.24.51 -> tools/gen1.py (auto); review=proofread-sonnet; flags=none
\begin{question}{77}
(주)감평은 분권화된 사업부 A와 B를 이익중심점으로 운영한다. 사업부 A는 매년 부품 X를 8,000단위 생산하여 전량 외부시장에 단위당 ￦150에 판매하여 왔다. 최근 사업부 B는 제품 단위당 부품 X가 1단위 소요되는 신제품 Y를 개발하고, 단위당 판매가격 ￦350에 4,000단위를 생산·판매하는 방안을 검토하고 있다. 다음은 부품 X에 대한 제조원가와 신제품 Y에 대한 예상제조원가 관련 자료이다.
\par{\small\centering\begin{tabular}{|c|c|c|}\hline
구 분 & 부품 X & 신제품 Y \\ \hline
단위당 직접재료원가 & ￦40 & ￦80 \\ \hline
단위당 직접노무원가 & 35 & 70 \\ \hline
단위당 변동제조간접원가 & 25 & 30 \\ \hline
연간 고정제조간접원가 & 200,000 & 100,000 \\ \hline
연간 최대생산능력 & 10,000단위 & 5,000단위 \\ \hline
\end{tabular}\par}
사업부 B는 신제품 Y의 생산에 필요한 부품 X를 사내대체하거나 외부로부터 단위당 ￦135에 공급받을 수 있다. 사업부 A는 사내대체를 전량 수락하든지 기각해야 하며, 사내대체 시 외부시장 판매를 일부 포기해야 한다. 사업부 A가 사내대체를 수락할 수 있는 부품 X의 단위당 최소대체가격은?
\begin{choices}
  \item ￦100
  \item ￦125
  \item ￦135
  \item ￦170
  \item ￦180
\end{choices}
\end{question}
\answer{2}
